\section{Mezclador de vídeo}

El mezclador de vídeo ---también denominado \textit{switcher}--- es el componente
central de cualquier sistema de producción audiovisual en directo. Su función es
seleccionar, en cada instante, qué fuente de vídeo (o combinación de fuentes) se emite
a la salida de programa, gestionando además las transiciones entre fuentes y la
incorporación de elementos gráficos sobre la imagen.

\subsection{Mezclado tradicional}

En la producción broadcast tradicional, el mezclado se realiza mediante hardware
dedicado. Los mezcladores hardware profesionales (Blackmagic ATEM, Ross Video,
Grass Valley Kayenne) ofrecen latencias sub-imagen, alta fiabilidad y soporte para
múltiples buses simultáneos de programa y preview. Son la solución estándar en
producciones televisivas de gran escala: retransmisiones deportivas, eventos
parlamentarios, galas de premios y producciones multicámara de cualquier tipo.

Un caso de uso paradigmático es la producción de retransmisiones en directo de grandes
eventos, donde un mezclador hardware gestiona decenas de fuentes de cámara simultáneas
con conmutación en tiempo real, garantizando la continuidad de la señal incluso ante
fallos puntuales. Esta robustez se logra mediante redundancia hardware (fuentes de
alimentación duplicadas, procesadores de backup) y certificación de los equipos para
operación continua.

\subsection{Mezclado por software y producción remota}

La producción audiovisual basada en software de mezclado permite realizar la producción
sin hardware dedicado, utilizando ordenadores de propósito general. Este enfoque ha
democratizado el acceso a la producción audiovisual profesional, reduciendo la barrera
de entrada tanto en coste como en complejidad de instalación.

Las principales herramientas de mezclado por software son:

\begin{itemize}
  \item \textbf{OBS Studio}: solución de código abierto ampliamente utilizada para
  \textit{streaming} y grabación. Su arquitectura monolítica está orientada a un único
  operador en un único equipo, lo que la hace ideal para producciones individuales pero
  limita su uso en entornos distribuidos donde el procesamiento debe separarse de la
  interfaz de operador.

  \item \textbf{vMix}: solución comercial profesional para Windows con soporte para
  NDI, multiview y producción 4K. Su coste de licencia y su limitación a Windows la
  excluyen de entornos de producción \textit{open-source} sobre Linux.

  \item \textbf{Voctomix}: herramienta de código abierto con arquitectura
  cliente-servidor que separa el motor de mezcla (\texttt{voctocore}) de la interfaz
  de control (\texttt{voctogui}). Esta separación permite que el procesamiento se
  realice en un servidor mientras el operador trabaja desde otro equipo conectado por
  red, facilitando la producción completamente remota. Es la herramienta desarrollada
  y extendida en el presente trabajo.
\end{itemize}

La investigación académica sobre mezclado de vídeo por software en entornos remotos
ha crecido en los últimos años de la mano de la madurez de las infraestructuras de
red IP. Trabajos como \cite{XXX} han analizado la viabilidad de arquitecturas
cliente-servidor para producción en directo sobre redes de campus universitarios,
concluyendo que herramientas como Voctomix son adecuadas para producciones de mediana
escala cuando se dispone de una red local de baja latencia.
% TODO: añadir referencias bibliográficas relevantes sobre mezclado remoto y producción IP

\subsection{GStreamer como motor de mezclado}

GStreamer es el \textit{framework} multimedia de código abierto sobre el que se
construye voctocore. Su arquitectura de \textit{pipeline} ---grafos dirigidos de
elementos de procesamiento interconectados mediante \textit{pads} tipados--- permite
construir cadenas de composición de vídeo arbitrariamente complejas combinando
elementos existentes. El elemento central en Voctomix es \texttt{compositor}, el
mezclador de vídeo multipista de GStreamer 1.x, que acepta $N$ pads de entrada con
sus propias transformaciones geométricas y produce un único flujo de salida. Las
propiedades de cada pad se modifican en tiempo real mediante mensajes GObject,
permitiendo transiciones suaves sin interrumpir el pipeline ni generar fotogramas
corruptos.
